\section{Background and Related Work}

The gap between what retrieval systems deliver and what scientific data discovery demands spans four areas: the hybrid retrieval foundations we build upon, the numeracy crisis that motivates explicit arithmetic operators, quantity-aware retrieval systems that attempt but fail to solve cross-unit matching, and scientific data discovery in HPC.

\subsection{Hybrid Retrieval Foundations}

Modern retrieval combines lexical and dense signals. BM25 \cite{robertson2009bm25} remains the strongest lexical baseline; dense retrieval via learned bi-encoders \cite{karpukhin2020dpr} captures semantic similarity that keyword matching misses. Bruch et al.\ \cite{bruch2023fusion} show that a learned convex combination consistently outperforms reciprocal rank fusion (RRF) on TREC benchmarks, establishing the analytic framework we adopt. At the agentic frontier, A-RAG \cite{arag2026} introduces hierarchical retrieval interfaces that let LLM agents compose retrieval strategies dynamically.

We build \emph{on} hybrid retrieval, not beside it. Our science-aware operators execute as parallel branches alongside BM25 and dense search within the same fusion pipeline, adding domain-specific recall without replacing any existing component.

\subsection{The Numerical Reasoning Crisis}

Numbers are fundamentally broken in both retrieval systems and large language models. The Numeracy Gap study \cite{numeracygap2026} evaluates 13 embedding models on numerical content and reports 0.54 binary retrieval accuracy---barely above random chance. NC-Retriever \cite{ncretriever2025} finds that dense retrieval achieves only 16.3\% accuracy on queries with numeric constraints, even when the constraint is stated explicitly in both query and document. The root cause is architectural: tokenizers encode numbers per-digit rather than by magnitude, so ``200'' and ``200000'' share no meaningful representation \cite{numbers_not_continuous2025}. GSM-Symbolic \cite{gsmsymbolic2025} delivers the sharpest result: changing only the numeric values in grade-school math problems while preserving problem structure causes a 65\% accuracy drop, proving that LLMs pattern-match on surface numerals rather than reason about magnitude.

This is not merely a unit problem. Numbers themselves---the substrate of all scientific measurement---are broken in every retrieval and language model. Scientific data retrieval, where matching requires comparing magnitudes across units, is especially vulnerable.

\subsection{Quantity-Aware Retrieval}

Three lineages attempt to make retrieval quantity-aware. None succeeds at arithmetic conversion.

\textbf{String normalization.} QFinder \cite{qfinder2022} extends Elasticsearch with quantity-aware ranking by extracting and normalizing numeric expressions. CQE \cite{cqe2023} advances this with dependency parsing over a 531-unit dictionary, handling complex quantity expressions. Numbers Matter! \cite{numbersmatter2024} builds a full retrieval pipeline with quantity-aware indexing and evaluation on FinQuant and MedQuant benchmarks. The fundamental limitation persists: string normalization cannot cross SI prefixes. ``kilopascal'' and ``pascal'' remain distinct strings after normalization, so a query for ``200~kPa'' never matches a document containing ``200000~Pa.''

\textbf{Learned embeddings.} CONE \cite{cone2026} trains a transformer to embed number--unit pairs into a shared vector space, achieving a 25\% Recall@10 gain on quantity-bearing queries. NC-Retriever \cite{ncretriever2025} applies contrastive learning to numeric constraint matching. These approaches are probabilistic: a model may learn that 200~kPa is \emph{similar} to 200000~Pa, but provides no correctness guarantee. Weller et al.\ \cite{embedding_limits2025} prove theoretical upper bounds on embedding-based retrieval, showing that fixed-dimensional embeddings cannot faithfully preserve all distance relationships---a result that undercuts learned approaches to exact quantity matching.

\textbf{Math-aware search.} Approach0 \cite{approach0} matches equation structure via operator trees. SSEmb \cite{ssemb2025} learns semantic similarity over mathematical expressions. These systems match formula \emph{structure} but are disconnected from physical measurements---they cannot link $F = ma$ to a document reporting force in newtons.

\textbf{Offline conversion tools.} Unit conversion itself is well-established: UDUNITS-2~\cite{udunits2} and GNU Units convert data values during preprocessing, but they do not operate at retrieval time, do not index measurements, and expose no search interface---a scientist must already know which files to convert.

\textbf{Gap.} To our knowledge, no existing \emph{retrieval system} integrates arithmetic SI conversion as a first-class retrieval operator. Offline tools require locating data first; string normalization and learned embeddings operate within retrieval but cannot perform arithmetic. No system combines quantity-aware retrieval with formula normalization or federated search across HPC backends.

\subsection{Scientific Data Discovery in HPC}

\textbf{Paper search.} OpenScholar \cite{openscholar2026} indexes 45 million open-access papers and achieves GPT-4o-level citation quality. HiPerRAG \cite{hiperrag2025} scales RAG to 3.6 million articles across DOE supercomputers Polaris, Sunspot, and Frontier. Both search \emph{papers}---not experimental data, simulation outputs, or HDF5 files.

\textbf{Data search.} Recent systems do search scientific data. PANGAEA-GPT \cite{pangaeagpt2026} deploys hierarchical multi-agent search over 400,000+ geoscientific datasets with data-type-aware routing and unit scale validation, achieving 8.14/10 relevance versus a 2.87/10 baseline. LLM-Find \cite{llmfind2025} uses LLM-generated code to retrieve geospatial data from heterogeneous sources with 80--90\% success. The distinction is temporal: PANGAEA-GPT validates unit scales \emph{post-hoc} (after retrieval); we convert \emph{at retrieval time} (during indexing and query processing). LLM-Find generates access code per source; we provide reusable retrieval operators.

\textbf{HPC data management.} Datum \cite{datum2024} catalogs scientific data at INL but offers no content search. BRINDEXER \cite{brindexer2020} indexes POSIX metadata on Lustre with 69\% query improvement---metadata only, no content. PROV-IO+ \cite{provio2024} captures HDF5 I/O provenance with under 3.5\% overhead but exposes no search interface. HDF Clinic \cite{hdfclinic2025} proposes a knowledge graph over HDF5 but remains unimplemented.

\textbf{Agentic retrieval.} Context-1~\cite{context1_2026} trains a 20B-token agentic search agent; PRISM~\cite{prism2025} decomposes retrieval into precision/recall subagents; Agentic-R~\cite{agenticr2026} co-trains retrievers with answer-correctness rewards (10--13\% multi-hop gains). CLADD~\cite{cladd2025} and SciAgent~\cite{sciagent2025} add domain knowledge operators for drug discovery and scientific workflows. These systems advance orchestration and learned retrieval. To our knowledge, none incorporates arithmetic dimensional conversion or formula normalization into the retrieval pipeline. Our work is complementary: PRISM and Agentic-R could be layered on top of our science-aware operators.

\begin{table*}[!t]
\centering
\caption{Capability comparison across retrieval and scientific data systems. \cmark~=~supported, \xmark~=~not supported, $\circ$~=~partial.}
\label{tab:comparison}
\small
\renewcommand{\arraystretch}{1.2}
\setlength{\tabcolsep}{5pt}
\begin{tabular}{l c c c c c c c c c c}
\toprule
 & \multicolumn{5}{c}{\textbf{Our Contributions}} & \multicolumn{5}{c}{\textbf{General Capabilities}} \\
\cmidrule(lr){2-6} \cmidrule(lr){7-11}
\textbf{System}
  & \rotatebox{60}{\small SI Conversion}
  & \rotatebox{60}{\small Formula Norm.}
  & \rotatebox{60}{\small Sci.\ Branches}
  & \rotatebox{60}{\small Federated}
  & \rotatebox{60}{\small HDF5/NetCDF}
  & \rotatebox{60}{\small Agentic}
  & \rotatebox{60}{\small Qty.\ Extract.}
  & \rotatebox{60}{\small Paper Search}
  & \rotatebox{60}{\small Data Search}
  & \rotatebox{60}{\small LLM Integr.} \\
\midrule
\rowcolor{gray!15}
\textbf{Ours (clio)}  & \cmark & \cmark & \cmark & \cmark & \cmark & \cmark & \cmark & \xmark & \cmark & \cmark \\
PANGAEA-GPT            & \xmark & \xmark & \xmark & \xmark & \xmark & \cmark & $\circ$ & \xmark & \cmark & \cmark \\
\rowcolor{gray!5}
OpenScholar            & \xmark & \xmark & \xmark & \xmark & \xmark & \xmark & \xmark & \cmark & \xmark & \cmark \\
HiPerRAG               & \xmark & \xmark & \xmark & \xmark & \xmark & \xmark & \xmark & \cmark & \xmark & \cmark \\
\rowcolor{gray!5}
Numbers Matter!        & \xmark & \xmark & \xmark & \xmark & \xmark & \xmark & \cmark & \xmark & \xmark & \xmark \\
CONE                   & \xmark & \xmark & \xmark & \xmark & \xmark & \xmark & \cmark & \xmark & \xmark & \xmark \\
\rowcolor{gray!5}
CLADD                  & \xmark & \xmark & \xmark & \xmark & \xmark & \cmark & \xmark & \xmark & \xmark & \cmark \\
SciAgent               & \xmark & \xmark & \xmark & \xmark & \xmark & \cmark & \xmark & \xmark & \xmark & \cmark \\
\rowcolor{gray!5}
Context-1              & \xmark & \xmark & \xmark & \xmark & \xmark & \cmark & \xmark & \xmark & \xmark & \cmark \\
\bottomrule
\end{tabular}
\renewcommand{\arraystretch}{1.15}
\end{table*}

Table~\ref{tab:comparison} summarizes the landscape. Individual capabilities exist in isolation---UDUNITS converts units offline, PANGAEA-GPT validates scales post-hoc, OpenScholar searches papers at scale---but to our knowledge, no prior \emph{retrieval system} integrates arithmetic SI conversion as a first-class retrieval operator, combines it with formula normalization, or operates across federated HPC backends with per-connector capability negotiation. We honestly mark capabilities we lack---we do not index papers. PANGAEA-GPT receives partial credit for quantity extraction (unit scale validation, but post-hoc rather than at retrieval time).

%% ===================================================================
%% Section 3: System Design
